% ============================================================
%  Příprava na maturitu – Mechatronika
%  Kompiluj: pdflatex maturita_mechatronika.tex
%  Kódování: UTF-8
% ============================================================

\documentclass[a4paper,10pt]{article}

% --- Balíčky ---
\usepackage[T1]{fontenc}
\usepackage[utf8]{inputenc}
\usepackage[czech]{babel}
\usepackage[left=2.4cm, right=1.8cm, top=1.8cm, bottom=1.8cm]{geometry}
\usepackage{lmodern}
\usepackage{microtype}
\usepackage{parskip}
\usepackage{titlesec}
\usepackage{enumitem}
\usepackage{xcolor}
\usepackage{hyperref}
\usepackage{tabularx}
\usepackage{booktabs}
\usepackage{array}
\usepackage{tikz}
\usepackage{eso-pic}
\usepackage{mdframed}
\usepackage{amsmath}
\usepackage{amssymb}
\usepackage{pgfplots}
\pgfplotsset{compat=1.18}
\usetikzlibrary{arrows.meta, positioning, shapes.geometric, calc}

% --- Barvy ---
\definecolor{accent}{HTML}{03254E}
\definecolor{stripeblue}{HTML}{568EA3}
\definecolor{medgray}{HTML}{888888}
\definecolor{lightblue}{HTML}{EAF2F8}
\definecolor{lightyellow}{HTML}{FDFBE4}
\definecolor{lightgreen}{HTML}{EAF7EA}

% --- Modrý pruh vlevo ---
\AddToShipoutPictureBG{%
  \begin{tikzpicture}[remember picture, overlay]
    \fill[stripeblue]
      (current page.north west)
      rectangle
      ([xshift=10mm] current page.south west);
    \fill[accent!60]
      ([xshift=10mm] current page.north west)
      rectangle
      ([xshift=12mm] current page.south west);
  \end{tikzpicture}%
}

\hypersetup{colorlinks=true, urlcolor=accent, linkcolor=accent}

\titleformat{\section}
  {\large\bfseries\color{accent}}
  {\thesection.}
  {0.5em}
  {}
  [\color{accent}\titlerule]

\titlespacing{\section}{0pt}{16pt}{6pt}

\newmdenv[
  backgroundcolor=lightblue,
  linecolor=stripeblue,
  linewidth=1pt,
  roundcorner=4pt,
  innerleftmargin=8pt,
  innerrightmargin=8pt,
  innertopmargin=6pt,
  innerbottommargin=6pt,
  skipabove=4pt,
  skipbelow=4pt
]{pojmybox}

\newmdenv[
  backgroundcolor=lightyellow,
  linecolor=accent!40,
  linewidth=1pt,
  roundcorner=4pt,
  innerleftmargin=8pt,
  innerrightmargin=8pt,
  innertopmargin=6pt,
  innerbottommargin=6pt,
  skipabove=4pt,
  skipbelow=8pt
]{vysvetlenibox}

\newmdenv[
  backgroundcolor=lightgreen,
  linecolor=accent!30,
  linewidth=1pt,
  roundcorner=4pt,
  innerleftmargin=8pt,
  innerrightmargin=8pt,
  innertopmargin=6pt,
  innerbottommargin=6pt,
  skipabove=4pt,
  skipbelow=8pt
]{vzorcebox}

\pagestyle{plain}

% ============================================================
\begin{document}

% --- Záhlaví ---
\begin{minipage}[t]{0.75\textwidth}
    {\Huge\bfseries\color{accent} Příprava na maturitu}\\[4pt]
    {\large\color{stripeblue} Mechatronika}\\[2pt]
    \textcolor{medgray}{\small Straka Jakub, Suchánek Lukáš \quad SPŠ Prosek \quad 2026}
\end{minipage}
\hfill
\begin{minipage}[t]{0.22\textwidth}
    \raggedleft
    \begin{tikzpicture}
        \draw[color=stripeblue, rounded corners=4pt, line width=1pt]
            (0,0) rectangle (3,1.2)
            node[midway, align=center, color=accent, font=\small\bfseries]
            {Otázek: 22};
    \end{tikzpicture}
\end{minipage}

\vspace{6pt}
\noindent\color{accent}\rule{\linewidth}{1.5pt}
\vspace{2pt}

% ============================================================
\section{Číslicově řízené stroje (CNC) a jejich programování}

\begin{pojmybox}
\textbf{Klíčové pojmy:}
G-kód, M-kód, souřadnicový systém, absolutní/přírůstkové programování,
nulový bod, korekce nástroje, frézování, soustružení, CAM,
interpolace, posuvová rychlost, otáčky vřetena
\end{pojmybox}

\begin{vysvetlenibox}
\textbf{Princip CNC stroje}\\
Existují primárně 3 druhy CNC strojů - 3 osy, 4 osy, 5 os, a jejich kombinace (např 3+1 znamený tří-osý se čtvrtou osou přidanou)
\textbf{Souřadnicový systém CNC stroje}\\
Pravotočivý kartézský systém (X, Y, Z + rotační osy A, B, C).
Nulový bod obrobku (WPZ) se nastavuje před obráběním. Nulový bod stroje (MZ) se nastaví hned po zapnutí stroje například koncovými spínači a bez jeho nastavení jede stroj na slepo, jelikož neví kde se nachází.

\textbf{Přehled G-kódů}

\begin{center}
\begin{tabularx}{0.95\linewidth}{@{} l X @{}}
    \toprule
    \textbf{G-kód} & \textbf{Funkce} \\
    \midrule
    G00 & Rychloposuv (bez řezu) \\
    G01 & Lineární interpolace (řez přímou čarou) \\
    G02 & Kruhová interpolace po směru hodinových ručiček \\
    G03 & Kruhová interpolace proti směru hod. ručiček \\
    G17/18/19 & Výběr roviny (XY / XZ / YZ) \\
    G90/G91 & Absolutní / přírůstkové souřadnice \\
    G94/G95 & Posuv [mm/min] / [mm/otáčku] \\
    \bottomrule
\end{tabularx}
\end{center}

\textbf{Přehled M-kódů}

\begin{center}
\begin{tabular}{@{} l l @{\quad} l l @{}}
    \toprule
    M00 & Pauza & M08 & Chlazení zapnout \\
    M03 & Vřeteno CW & M09 & Chlazení vypnout \\
    M04 & Vřeteno CCW & M30 & Konec programu \\
    M05 & Vřeteno stop & M06 & Výměna nástroje \\
    \bottomrule
\end{tabular}
\end{center}

\textbf{Ukázkový blok G-kódu:}
\begin{vzorcebox}
\begin{verbatim}
N10 G90 G17 G94         ; Absolutní, rovina XY, mm/min
N20 T01 M06             ; Nástroj č. 1
N30 S1200 M03           ; 1200 RPM, vřeteno CW
N40 G00 X0 Y0 Z5        ; Rychloposuv nad nulový bod
N50 G01 Z-3 F100        ; Hloubkový záběr 3mm, F=100
N60 G01 X50 Y0 F300     ; Frézování přímkou X+50
N70 G02 X50 Y50 I0 J25  ; Oblouk CW, střed I0 J25
N80 G00 Z50             ; Zdvih nástroje
N90 M05 M30             ; Stop vřetena, konec
\end{verbatim}
\end{vzorcebox}

\textbf{CAM -- Computer Aided Manufacturing}\\
CAM software (Fusion 360, Mastercam, NX, SolidCAM, FreeCAD CAM) generuje G-kód z nakreslených/navrhlých cest. Postprocesor převede neutrální CLData na G-kód specifický
pro daný řídicí systém (Fanuc, Siemens Sinumerik, Heidenhain, GRBL, Hase).

\begin{vzorcebox}
Řezná rychlost: $v_c = \frac{\pi \cdot d \cdot n}{1000}$ [m/min]\\
$\Rightarrow$ Otáčky: $n = \frac{1000 \cdot v_c}{\pi \cdot d}$ [min$^{-1}$]\\
Minutový posuv: $v_f = f_z \cdot z \cdot n$ [mm/min]
\quad ($f_z$ = posuv na zub, $z$ = počet zubů)
\end{vzorcebox}
\end{vysvetlenibox}

% ============================================================

% ============================================================
\end{document}
